\section{图解：极限的误差带}
\begin{figure}[htbp]\centering
\begin{tikzpicture}\begin{axis}[width=.8\textwidth,height=6cm,axis lines=middle,xmin=0,xmax=1.8,ymin=0,ymax=2.9,xlabel=$x$,ylabel=$f(x)$,samples=100]
\addplot[Teal,thick,domain=0:1.65]{x^2};
\addplot[dashed,Gold,domain=0:1.7]{1.3};\addplot[dashed,Gold,domain=0:1.7]{.7};
\draw[dashed,Ink] (axis cs:.9,0)--(axis cs:.9,2.5);
\draw[dashed,Ink] (axis cs:1.1,0)--(axis cs:1.1,2.5);
\fill[Ink] (axis cs:1,1) circle[radius=1.8pt];
\node[anchor=west] at (axis cs:1.25,1.35) {$L+\varepsilon$};
\node[anchor=west] at (axis cs:1.25,.65) {$L-\varepsilon$};
\node at (axis cs:1,2.65) {$|x-1|<\delta$};
\end{axis}\end{tikzpicture}
\caption{以$f(x)=x^2,x\to1$为例，竖直带控制自变量，水平带控制函数误差。图取$\delta=0.1,\varepsilon=0.3$；当$|x-1|<0.1$时$|x^2-1|=|x-1||x+1|<0.21<0.3$。定义要求对任意正$\varepsilon$都能选$\delta$，图中只展示其中一次选择。}\label{fig:review:limit}
\end{figure}
